% =====================================================================
% Abstract (kurz, ca. 150-180 Wörter)
% Studienarbeit: RAG vs. Fine-Tuning - Ein konzeptioneller Kriterienkatalog
% Autor: Tobias Weigold
% =====================================================================

\begin{abstract}
Große Sprachmodelle stehen im betrieblichen Wissensmanagement in Spannung zu ihren strukturellen Grenzen wie Halluzinationen, statischem Wissensstand und fehlender Domänenexpertise. Zur Adressierung dieser Defizite haben sich Retrieval-Augmented Generation (RAG) und Fine-Tuning als dominante Anpassungsverfahren etabliert. Bisherige Vergleichsstudien fokussieren jedoch überwiegend auf technische Einzelaspekte und lassen organisatorische, wirtschaftliche und insbesondere regulatorische Dimensionen weitgehend unberücksichtigt. Die vorliegende Arbeit entwickelt einen konzeptionellen Kriterienkatalog aus fünfzehn Kriterien, der wissens-, aufgaben-, ressourcen- und regulatorische Dimensionen zusammenführt und explizit Anforderungen aus der DSGVO und dem EU AI Act berücksichtigt. Methodisch handelt es sich um eine theoretisch abgeleitete konzeptionelle Analyse auf Basis der Fachliteratur. Die Kriterien werden zu einem zweistufigen Entscheidungsmodell verdichtet, in dem regulatorische Anforderungen als K.-o.-Filter wirken und die übrigen Kriterien anschließend in einer Ordinalmatrix bewertet werden. Ein Fünf-Schritte-Leitfaden und drei illustrative Anwendungsszenarien machen das Modell praktisch anwendbar und zeigen, dass RAG in dynamischen, regulierten Kontexten strukturell im Vorteil ist, während Fine-Tuning bei stabilen, stilistisch geprägten Aufgaben und hybride Ansätze bei fachsprachlich anspruchsvollen Domänen tragfähig bleiben.
\end{abstract}
